\section{Limitations and Future Work}
\label{sec:limitations}


% ══════════════════════════════════════════════════════════════
\subsection{Limitations}
\label{sec:limits}

We identify six limitations that constrain the generalizability and statistical strength of the findings presented in this paper.

\paragraph{Single framework.}
The Completion Integrity System is implemented in the VIBE Framework for Claude Code.
All mechanical enforcement (Layers~1--3) is tested in this environment.
Generalizability to other agentic frameworks---Cursor, GitHub Copilot Workspace, Devin, Windsurf---is not demonstrated.
The taxonomy and root cause analysis draw on evidence from multiple systems, but the mitigation architecture has been validated in one.

\paragraph{Limited portability.}
Layers~1 and~2 of the Completion Integrity System depend on the Claude Code hook system, which provides shell-level interception of assistant messages before delivery.
Other CLI tools lack equivalent infrastructure.
Cross-model testing in the experimental evaluation uses prompt-only mitigation for non-Claude models because the mechanical enforcement layer cannot be ported without a compatible hook API.
This means the full-system condition was tested only on Claude, while the prompt-only and baseline conditions provide cross-model comparisons.

\paragraph{Small experimental scale.}
The experimental evaluation uses 20 tasks across 3 models, yielding 180 total experimental runs across 3 conditions.
This provides initial evidence for the central claims but lacks statistical power for fine-grained comparisons between models or between task categories.
Effect sizes for the mechanical vs.\ prompt-only comparison are [to be updated with experimental results], but confidence intervals are wide.
A larger-scale replication is needed to establish precise effect magnitudes and to test for interaction effects between model, task type, and mitigation strategy.

\paragraph{Manual evaluation.}
Completeness scoring in the experimental evaluation requires human judgment.
A human evaluator examines each model output against a predefined rubric of required items and assigns a completeness score.
Inter-rater reliability was not tested: a single evaluator scored all outputs.
Subjectivity in distinguishing ``found'' (item genuinely analyzed with evidence) from ``partial'' (item mentioned but not substantively analyzed) may affect individual scores.
The binary classification of false completion events (score below threshold) is less sensitive to this subjectivity than the continuous completeness scores, but the limitation remains.

\paragraph{Single point in time.}
The experimental results reflect specific model versions at testing time.
Model providers update their models regularly, and the false completion behavior documented here may change---improving, worsening, or shifting to different failure patterns---with subsequent releases.
The taxonomy provides a framework for tracking these changes, but the quantitative results are a snapshot, not a trend.
Longitudinal tracking is needed to determine whether the problem is stable, improving, or deteriorating.

\paragraph{Economic estimates.}
The industry-wide cost projections presented in \S\ref{sec:economic} rely on extrapolation from limited data points: the AMD dataset (one organization, one tool, one engineering domain) and community reports (self-selected, non-random sample).
The 80$\times$ cost inflation figure reflects a worst-case degradation event, not steady-state conditions.
Confidence intervals on the aggregate economic estimates are wide, and the projections should be treated as order-of-magnitude indicators rather than precise forecasts.


% ══════════════════════════════════════════════════════════════
\subsection{Future Work}
\label{sec:future-work}

Six directions extend the work presented in this paper.

\paragraph{Production telemetry.}
Deploy sentinel logging at scale in production environments.
Collect interception rates (how often the sentinel blocks a false completion), false positive rates (how often legitimate completions are blocked), and user satisfaction data over months of use.
Production telemetry would replace the controlled experimental estimates with ecologically valid measurements and enable calibration of detection thresholds.

\paragraph{Cross-framework validation.}
Port the sentinel concept to other agentic tools: Qwen CLI, Gemini CLI, and other systems that provide hook or middleware infrastructure.
Test whether mechanical enforcement generalizes across model families and tool architectures.
Identify the minimum infrastructure requirements for effective false completion detection.

\paragraph{Automated evaluation.}
Train a classifier to score task completeness from model outputs, reducing reliance on manual evaluation and enabling experimental protocols at larger scale.
An automated evaluator would allow testing across hundreds of tasks and model versions, providing the statistical power that the current 20-task protocol lacks.

\paragraph{Longitudinal tracking.}
Monitor false completion rates across model versions to test whether the problem improves, worsens, or remains stable over time.
The inverse scaling hypothesis (\S\ref{sec:disc-scaling}) predicts that false completion will worsen with model scale; longitudinal data would confirm or refute this prediction.

\paragraph{Community replication.}
Publish the experimental task definitions, scoring rubrics, and analysis scripts for independent validation.
The reproducibility of the findings depends on whether other researchers, using different evaluators and environments, observe the same patterns.

\paragraph{Formal verification.}
Explore whether the VIBE\_GATE marker protocol can be extended to provide cryptographic guarantees of tool execution.
If each tool invocation produces a signed attestation---a hash of the input, output, and timestamp---then the sentinel could verify not just that a tool was called, but that the recorded output is authentic and unmodified.
This would close the remaining trust gap between transcript evidence and ground truth.
